% =============================================================
% Capitolo: Strategie di Prompting
% Progetto: MCP_ECOM — Agente e-commerce basato su MCP
% =============================================================

\chapter{Strategie di Prompting}
\label{chap:prompting}

% ─────────────────────────────────────────────────────────────
\section{Introduzione}
\label{sec:prompt-intro}
% ─────────────────────────────────────────────────────────────

Nel progetto MCP\_ECOM il prompting non è un dettaglio implementativo ma
un vero e proprio \emph{layer architetturale} progettato esplicitamente.
Ogni chiamata all'LLM è governata da istruzioni costruite in modo
programmatico: quale ruolo assume il modello, che formato deve rispettare,
quante informazioni contestuali riceve e con quale priorità vengono
composte.

Questo approccio si contrappone al prompting implicito — dove il testo
di istruzione è disseminato casualmente nel codice — e introduce invece
un \textbf{sistema di strati a priorità crescente}:

\begin{enumerate}
    \item \textbf{Istruzioni personalizzate dell'utente} (priorità
        assoluta) — regole esplicite impostate a runtime che sovrascrivono
        qualsiasi comportamento di default.
    \item \textbf{Tono conversazionale} (alta priorità) — modifica lo
        stile senza alterare la logica.
    \item \textbf{System prompt di base} (priorità normale) — definisce
        persona, politiche e schema di output.
    \item \textbf{Contesto e memoria} (priorità contestuale) — scratchpad,
        storico di sessione, preferenze a lungo termine.
\end{enumerate}

Il file centrale che contiene le costanti e le funzioni di composizione è
\texttt{app/agent/prompts.py}; ulteriori prompt specializzati si trovano
distribuiti nei sottomoduli
(\texttt{app/services/parser.py},
\texttt{app/services/rag/query\_expansion.py},
\texttt{app/services/nlp\_ner.py},
\texttt{app/services/rag/ltr\_trainer.py},
\texttt{app/services/nlp\_sentiment.py},
\texttt{app/mcp/tools/conversation.py},
\texttt{app/llm/vision.py}).


% ─────────────────────────────────────────────────────────────
\section{Architettura del Prompting}
\label{sec:prompt-arch}
% ─────────────────────────────────────────────────────────────

\subsection{I Quattro System Prompt Principali}
\label{subsec:four-prompts}

Il sistema prevede quattro system prompt distinti, ciascuno attivo in una
fase diversa del ciclo di vita della richiesta:

\begin{description}
    \item[\texttt{PLANNER\_SYSTEM\_PROMPT}] Governa il loop ReAct: decide
        quale tool invocare al passo successivo. È attivo per ogni
        iterazione del ciclo plan\,→\,execute.

    \item[\texttt{PLAYWRIGHT\_PLANNER\_SYSTEM\_PROMPT}] Variante
        specializzata del Planner attivata quando l'agente opera in
        modalità browser reale (Playwright). Sostituisce completamente
        il precedente per evitare contraddizioni tra strumenti API e
        strumenti browser.

    \item[\texttt{FINAL\_ANSWER\_SYSTEM\_PROMPT}] Usato nella fase
        finale, dopo che tutti i tool hanno prodotto i loro risultati.
        Genera la risposta definitiva in linguaggio naturale strutturato.

    \item[\texttt{CONVERSATION\_ANSWER\_SYSTEM\_PROMPT}] Risponde a
        domande puramente conversazionali (saluti, curiosità generali)
        senza invocare template di analisi.
\end{description}

\subsection{Funzioni di Composizione}
\label{subsec:builder-functions}

Le due funzioni \texttt{build\_planner\_prompt()} e
\texttt{build\_final\_answer\_prompt()} costruiscono dinamicamente il
messaggio completo da inviare all'LLM. Entrambe accettano parametri
come \texttt{custom\_instructions}, \texttt{tone}, \texttt{mcp\_mode} e
lo \texttt{scratchpad} corrente.

La struttura generale del prompt prodotto da
\texttt{build\_planner\_prompt()} è la seguente:

\begin{lstlisting}[language=,caption={Struttura logica del prompt del planner},label=lst:planner-structure]
[ISTRUZIONI PERSONALIZZATE - PRIORITA' MASSIMA]   <- se presenti
[TONO RICHIESTO: AMICHEVOLE/PROFESSIONALE]        <- se specificato
[SYSTEM PROMPT BASE]                              <- PLANNER o PLAYWRIGHT
Step: N/MAX
Available tools: { ... catalogo MCP compatto ... }
Vision description: ...                           <- se immagine allegata
User query: ...
Scratchpad: { ... stato corrente agent ... }
Return JSON only.
\end{lstlisting}

\subsection{Token Budgeting e Troncamento}
\label{subsec:token-budget}

Per proteggere la finestra di contesto dell'LLM, il sistema applica un
budget di \textbf{6000 token} (\texttt{MAX\_LLM\_PROMPT\_TOKENS}),
con una stima grezza di 4 caratteri per token. Il budget residuo viene
assegnato allo scratchpad: se supera il limite, i campi più pesanti
vengono troncati secondo una priorità fissa:

\begin{enumerate}
    \item \texttt{results} (lista prodotti ricercati)
    \item \texttt{feedbacks} (feedback venditore)
    \item \texttt{rag\_context} (chunk RAG recuperati)
    \item \texttt{recent\_tool\_results}
\end{enumerate}

Per ciascuno di questi campi, la lista viene ridotta ai primi 3 elementi
e viene aggiunto un flag \texttt{\_truncated: true} per segnalare al
modello che il dato è stato compresso.


% ─────────────────────────────────────────────────────────────
\section{Role Prompting — Le Quattro Personalità di ebayGPT}
\label{sec:role-prompting}
% ─────────────────────────────────────────────────────────────

Il progetto adotta sistematicamente il \textbf{role prompting}: ogni
system prompt assegna al modello un'identità precisa con caratteristiche,
vincoli e politiche operative proprie.

\subsection{Il Planner Strategico}
\label{subsec:planner-prompt}

Il \texttt{PLANNER\_SYSTEM\_PROMPT} definisce l'LLM come
\emph{``il cervello strategico di un assistente e-commerce avanzato''}
con il compito di decidere la \emph{prossima azione migliore}. Il prompt
contiene 12 regole di pianificazione numerate e una sezione dedicata
alla \emph{policy degli strumenti}:

\begin{lstlisting}[language=Python,caption={Estratto del PLANNER\_SYSTEM\_PROMPT (app/agent/prompts.py)},label=lst:planner-prompt]
PLANNER_SYSTEM_PROMPT = """
Sei ebayGPT, il cervello strategico di un assistente e-commerce
avanzato. Il tuo compito e' decidere la PROSSIMA AZIONE MIGLIORE.

REGOLE DI PIANIFICAZIONE:
1. Priorita' alla Ricerca: Se l'utente esprime un interesse per un
   prodotto, marca o categoria, DEVI usare `search_products`.
2. Gestione del Contesto (CRITICO): Se l'utente usa pronomi o
   riferimenti impliciti (es: "ne", "quello", "neri?"), DEVI guardare
   `session_memory` per ricostruire una query COMPLETA.
3. Pensiero in Italiano: Tutti i tuoi ragionamenti nel campo `thought`
   devono essere in ITALIANO.
7. JSON Rigido: Rispondi SOLO con un JSON valido.
8. Confronto Efficiente (CRITICO): Se l'utente chiede un confronto
   tra prodotti gia' trovati nei top_results, NON chiamare
   search_products. Usa direttamente `comparison`.
...

POLICY STRUMENTI:
- search_products: Discovery, shopping, prezzi.
- analyze_seller: Affidabilita', reputazione, trust.
- get_item_details: Specifiche tecniche di un item_id gia' trovato.
- market_trends: Prezzi medi sul web e trend (Google Trends + Shopping).
- ebay_scrape: OBBLIGATORIO se l'utente menziona 'Playwright' o
  'browser reale'.
- contact_seller: Genera link per contattare un venditore eBay.
- conversation: SOLO per chiacchiericcio senza intento di acquisto.
"""
\end{lstlisting}

Le 12 regole coprono scenari critici come la \emph{risoluzione di
pronomi} (regola~2), la \emph{ricerca proattiva} prima del verdetto
(regola~5), il \emph{confronto efficiente} senza ricerche ridondanti
(regola~8) e l'identificazione del venditore target (regola~12).

\subsection{La Modalità Browser Playwright}
\label{subsec:playwright-planner}

Quando l'agente opera con il browser Chromium reale, il system prompt
cambia completamente. Il \texttt{PLAYWRIGHT\_PLANNER\_SYSTEM\_PROMPT}
introduce regole specifiche per la navigazione passo-passo e la
gestione dei banner GDPR:

\begin{lstlisting}[language=Python,caption={Estratto del PLAYWRIGHT\_PLANNER\_SYSTEM\_PROMPT (app/agent/prompts.py)},label=lst:playwright-prompt]
"""
Sei ebayGPT in modalita' [BROWSER REALE PLAYWRIGHT]. Hai accesso a
un browser Chromium REALE, VISIBILE e PERSISTENTE.

REGOLE DI PIANIFICAZIONE INTERATTIVA:
1. Navigazione Step-by-Step (CRITICO): Non esistono macro-ricerche
   (NO search_products). Usa SEMPRE COME PRIMA AZIONE `browser_navigate`
   passando 'https://www.ebay.it'.
2. Attenzione ai Cookie Banner (CRITICO): E' OBBLIGATORIO fare
   browser_click su '#gdpr-banner-accept' per chiudere il popup.
3. Per cercare: usa browser_type passando il selettore '#gh-ac',
   la keyword, e impostando press_enter=true.
4. Analisi Visiva Costante: ogni step restituisce URL attuale, titolo
   e estratto di testo. Leggi le observations per capire dove sei!
6. Chiusura e Fallimento: se la ricerca restituisce 0 risultati,
   termina con 'action': 'finish' e spiega in final_answer.
"""
\end{lstlisting}

La separazione netta tra i due planner evita che l'LLM mescoli strumenti
API (come \texttt{search\_products}) con strumenti browser: un errore
che porterebbe a loop infiniti o chiamate invalide.

\subsection{Il Consulente Esperto — Final Answer}
\label{subsec:final-answer-prompt}

La fase di sintesi finale usa un prompt che trasforma l'LLM in un
\emph{personal shopper consulenziale}. La struttura della risposta è
imposta tramite un template a tre sezioni:

\begin{lstlisting}[language=Python,caption={Struttura template del FINAL\_ANSWER\_SYSTEM\_PROMPT (app/agent/prompts.py)},label=lst:final-answer-prompt]
FINAL_ANSWER_SYSTEM_PROMPT = """
Sei ebayGPT, il consulente esperto di shopping ufficiale. Il tuo
obiettivo e' guidare l'utente verso l'acquisto migliore.

STRUTTURA DELLA RISPOSTA (SEGUI QUESTO TEMPLATE):
## Analisi

[Analisi del contesto e tabella prodotti]

## Affidabilita'

[Analisi del trust score e dei venditori]

## Verdetto

[Consiglio finale]

REGOLE DI FORMATTAZIONE (CRITICO):
- DOPPIO INVIO: Dopo ogni titolo ## devi inserire DUE INVIO.
- FILTRI NEGATIVI (CRITICO): Se l'utente specifica di NON volere
  qualcosa (es. "no nero", "senza zip"), i prodotti DEVONO rispettare
  rigorosamente queste esclusioni.
- TABELLE: una sola tabella completa. La colonna Link deve essere l'ultima.
- URL OBBLIGATORI: non inventare mai link. Usa solo i dati forniti.
{TONE_PLACEHOLDER}
"""
\end{lstlisting}

Degna di nota è la regola sui \textbf{filtri negativi}: è esplicitamente
vietato al modello proporre articoli con attributi che l'utente ha
escluso, anche se i risultati di ricerca li contengono.

\subsection{Il Compagno di Chat}
\label{subsec:conversation-prompt}

Per le richieste puramente conversazionali (saluti, domande generiche
sull'e-commerce) viene usato un prompt minimale che enfatizza il tono
informale e la brevità:

\begin{lstlisting}[language=Python,caption={CONVERSATION\_ANSWER\_SYSTEM\_PROMPT (app/agent/prompts.py)},label=lst:conversation-prompt]
CONVERSATION_ANSWER_SYSTEM_PROMPT = """
Sei ebayGPT, un assistente e-commerce amichevole e competente.
Rispondi in modo naturale, colloquiale e conciso.

REGOLE:
- Tono caldo, diretto e personale. Usa "tu" informale.
- Per saluti o domande generali: 1-3 frasi, no elenchi, no titoli.
- MAI "Come posso aiutarti oggi?" come unica risposta: aggiungi
  sempre valore.
- Se hai storico di conversazione, usalo per risposte contestualizzate.
"""
\end{lstlisting}

Questo prompt è utilizzato sia direttamente nell'agente sia come base
nella funzione \texttt{\_build\_conversation\_prompt()} del tool MCP
\texttt{conversation}, che vi aggiunge dinamicamente lo storico della
sessione.


% ─────────────────────────────────────────────────────────────
\section{Chain-of-Thought e Output Strutturato}
\label{sec:cot-structured}
% ─────────────────────────────────────────────────────────────

\subsection{Il Campo \texttt{thought} come Ragionamento Esplicito}
\label{subsec:thought-field}

Il planner è istruito a produrre un campo \texttt{thought} in italiano
prima di scegliere l'azione. Questo è un'implementazione diretta del
\textbf{chain-of-thought prompting}: forzare il modello a esternalizzare
il ragionamento intermedio migliora la qualità della decisione finale,
in particolare per query ambigue.

Lo schema JSON di output imposto al planner è il seguente:

\begin{lstlisting}[language=,caption={Schema di output JSON del planner},label=lst:planner-schema]
{
  "thought":       "Spiega la strategia in ITALIANO",
  "intent":        "conversation | seller_analysis | product_search |
                    hybrid | comparison | item_details | shipping |
                    market_trends | deals | wishlist |
                    contact_seller | playwright_search",
  "action":        "tool_name | finish",
  "action_input":  { ... },
  "final_answer":  null
}
\end{lstlisting}

Il campo \texttt{intent} funge da classificatore di intenzione esplicito:
l'agente deve dichiarare l'intent prima di scegliere il tool, rendendo
il ragionamento verificabile.

\subsection{Strategia di Estrazione JSON}
\label{subsec:json-extraction}

L'output dell'LLM può contenere testo attorno al JSON (preamble,
markdown, blocchi \texttt{```json}). Il codice in
\texttt{app/agent/planner.py} applica una strategia di estrazione
a cascata:

\begin{enumerate}
    \item Cerca un blocco \texttt{```json ... ```} e ne estrae il contenuto.
    \item In assenza, invoca \texttt{extract\_first\_json\_object()} che
        individua le parentesi graffe bilanciate nel testo.
    \item Verifica che il tool indicato nel campo \texttt{action} esista
        nel catalogo MCP live (guard against hallucination).
    \item In caso di fallimento, genera un'azione di fallback
        (\texttt{finish}) con un messaggio di errore.
\end{enumerate}

\subsection{LTR Judge — Rilevanza con Vincoli}
\label{subsec:ltr-judge}

Un caso speciale di output strutturato è il \emph{giudice di rilevanza}
usato nella pipeline Learning-to-Rank. Il prompt istruisce l'LLM a
valutare ogni prodotto su una scala 0–4, con una rubrica dettagliata
che include il rispetto esplicito dei vincoli (prezzo, marca):

\begin{lstlisting}[language=Python,caption={JUDGE\_SYSTEM\_PROMPT per LTR (app/services/rag/ltr\_trainer.py)},label=lst:judge-prompt]
JUDGE_SYSTEM_PROMPT = """
Sei un esperto di e-commerce e Information Retrieval.
Valuta la rilevanza dei prodotti rispetto a una query con vincoli.

Scala 0-4:
4: Risultato perfetto. Rispetta TUTTI i vincoli (prezzo, marca).
3: Molto buono. Stessa categoria/marca, dettagli leggermente diversi.
2: Mediocre. Stessa categoria ma modello diverso, o accessori correlati.
1: Poco utile. Molto lontano dai vincoli o categoria vagamente correlata.
0: Irrilevante o rumore. Categoria completamente diversa.

IMPORTANTE: Prodotti che violano palesemente un vincolo di prezzo
o marca NON possono avere voto superiore a 1.
Rispondi solo con un JSON array: [4, 0, 1, 2]
"""
\end{lstlisting}

La regola critica che limita a 1 il voto massimo per le violazioni
di vincoli garantisce che il modello non sovrastime prodotti
che non rispettano i requisiti esplicitati dall'utente.


% ─────────────────────────────────────────────────────────────
\section{Tone Control Dinamico}
\label{sec:tone-control}
% ─────────────────────────────────────────────────────────────

Il sistema supporta tre modalità di tono conversazionale — \emph{amichevole},
\emph{professionale} e \emph{neutral} — che vengono iniettate nel prompt
a runtime tramite un meccanismo di \textbf{placeholder} e \textbf{prefisso}.

\subsection{Il Meccanismo \texttt{TONE\_PLACEHOLDER}}
\label{subsec:tone-placeholder}

Il \texttt{FINAL\_ANSWER\_SYSTEM\_PROMPT} contiene il token
\texttt{\{TONE\_PLACEHOLDER\}} nella sezione \emph{Tono e Persona}.
La funzione \texttt{build\_final\_answer\_prompt()} lo sostituisce
a runtime con la descrizione corrispondente al tono richiesto:

\begin{lstlisting}[language=Python,caption={Mapping tono → descrizione in build\_final\_answer\_prompt() (app/agent/prompts.py)},label=lst:tone-mapping]
tone_desc = {
    "amichevole":  "Usa un tono molto amichevole, caloroso e informale."
                   " Usa emoji e sii entusiasta.",
    "professionale": "Usa un tono formale, asciutto e professionale."
                     " Evita troppi fronzoli e sii preciso.",
    "neutral":     "Usa un tono equilibrato, educato e naturale."
}.get(tone.lower(), "Usa un tono naturale.")

system_prompt = system_prompt.replace("{TONE_PLACEHOLDER}",
                                      f"- {tone_desc}")
\end{lstlisting}

\subsection{Iniezione nel Planner}
\label{subsec:tone-planner}

Per il planner il tono viene iniettato come \textbf{prefisso} sopra il
system prompt base, prima ancora delle regole di pianificazione:

\begin{lstlisting}[language=Python,caption={Iniezione del tono come prefisso (app/agent/prompts.py)},label=lst:tone-prefix]
if tone:
    tone_instruction = (
        f"\nTONO RICHIESTO: {tone.upper()}. "
        "Mantieni questo stile in tutte le tue interazioni.\n"
    )
    system_prompt = tone_instruction + system_prompt
\end{lstlisting}

\subsection{Override tramite Custom Instructions}
\label{subsec:tone-override}

Le istruzioni personalizzate dell'utente vengono posizionate \emph{prima}
del tono e del system prompt base, garantendo che abbiano la precedenza
assoluta. Se l'utente specifica ``rispondi sempre in inglese'' nelle
custom instructions, questa regola sovrascrive anche il tono
predefinito in italiano.


% ─────────────────────────────────────────────────────────────
\section{Context Injection e Memoria Utente}
\label{sec:context-injection}
% ─────────────────────────────────────────────────────────────

Una delle strategie più articolate del progetto è l'iniezione contestuale
multipla: il prompt finale non contiene solo le istruzioni statiche ma
incorpora dinamicamente il profilo dell'utente, la storia della sessione
e le preferenze a lungo termine.

\subsection{Profilo Utente via Risorsa MCP}
\label{subsec:mcp-profile}

Prima di costruire il prompt, l'agente recupera il profilo utente tramite
una risorsa MCP (\texttt{profile://\{user\_id\}}). I dati estratti —
lingua preferita, brand favoriti, fascia di prezzo — vengono iniettati
nelle \texttt{custom\_instructions} passate al planner:

\begin{lstlisting}[language=Python,caption={Recupero ed iniezione del profilo MCP (app/agent/ebay\_agent.py)},label=lst:mcp-profile]
# Fetch dynamic user profile from MCP resource
profile_resource = await mcp_client.read_resource(
    f"profile://{user_id}"
)
if profile_resource:
    profile_data = json.loads(profile_resource)
    custom_instructions = (
        f"Lingua preferita: {profile_data.get('language', 'it')}\n"
        f"Brand preferiti: {profile_data.get('favorite_brands', [])}\n"
        f"Fascia di prezzo: {profile_data.get('price_preference', 'N/A')}"
    )
\end{lstlisting}

\subsection{Preferenze a Lungo Termine nel Final Answer}
\label{subsec:ltm-preferences}

La funzione \texttt{build\_final\_answer\_prompt()} inietta nel prompt
un blocco \texttt{USER STATUS/PREFERENCES} costruito dalla
\texttt{long\_term\_memory} dello scratchpad. Le preferenze estratte
includono:

\begin{itemize}
    \item Brand preferiti confermati (\texttt{db\_favorite\_brands})
    \item Brand cercati di recente (\texttt{validated\_brands})
    \item Budget medio abituale (\texttt{db\_price\_preference})
    \item Condizione preferita: nuovo / usato / ricondizionato
    \item Categorie di interesse (\texttt{category\_history})
    \item Profilo comportamentale: \emph{researcher} o \emph{power\_buyer}
    \item Venditori preferiti (\texttt{favorite\_sellers})
    \item Ricerche precedenti (\texttt{previous\_searches})
\end{itemize}

Questo blocco viene aggiunto come sezione esplicita del prompt con
l'intestazione \texttt{USER STATUS/PREFERENCES}, consentendo all'LLM
di personalizzare il verdetto finale in base alla storia dell'utente.

\subsection{Risoluzione dei Pronomi via Context Injection}
\label{subsec:pronoun-resolution}

Il parser semantico (\texttt{app/services/parser.py}) riceve un blocco
\texttt{CONVERSATION CONTEXT} che contiene i topic recenti della sessione.
Una regola esplicita nel prompt istruisce il modello a risolvere pronomi
impliciti come ``li'', ``quelli'', ``cercalo'' usando questo contesto:

\begin{lstlisting}[language=Python,caption={Iniezione del contesto conversazionale nel parser (app/services/parser.py)},label=lst:context-injection]
if context_info:
    context_block = (
        "\nCONVERSATION CONTEXT "
        "(History of recent topics/products):\n"
        f"{context_info}\n"
    )

# Nel prompt:
"IMPORTANT: Resolve pronouns (e.g., 'li', 'le', 'quelli',
 'cercalo', 'them', 'it') using the CONVERSATION CONTEXT."
\end{lstlisting}


% ─────────────────────────────────────────────────────────────
\section{Semantic Query Parsing}
\label{sec:query-parsing}
% ─────────────────────────────────────────────────────────────

Il modulo \texttt{app/services/parser.py} implementa un parser semantico
LLM-based che trasforma la query in linguaggio naturale dell'utente in
una struttura JSON normalizzata. Il prompt adotta il ruolo di
\emph{``strict semantic query parser''} e definisce uno schema di output
preciso con la distinzione tra \textbf{constraints} (requisiti obbligatori)
e \textbf{preferences} (desideri opzionali).

\subsection{Schema di Output e Tipi di Vincolo}
\label{subsec:parser-schema}

\begin{lstlisting}[language=Python,caption={Prompt del parser semantico (app/services/parser.py)},label=lst:parser-prompt]
prompt = f"""
You are a strict semantic query parser for an e-commerce assistant.
{context_block}
Return ONLY valid minified JSON. No markdown. No code fences.

Schema:
{{
  "semantic_query": "",
  "product": null,
  "product_type": "Main" | "Accessory" | "Part" | "Service" | "Unknown",
  "brands": [],
  "compatibilities": {{}},
  "constraints": [],
  "preferences": [],
  "target_seller": null
}}

Allowed constraint/preference types:

Price: {{"type":"price","operator":"<=" | ">=" | "between",
         "value": number OR [min,max]}}
Condition: {{"type":"condition","value":"new"|"used"|"refurbished"}}
Aspect: {{"type":"aspect","name":"Color"|"Storage"|"Size",
          "value":"Blue"|"256GB"|"42"}}

Rules:
- Put ONLY mandatory requirements in constraints.
- Put optional wishes in preferences.
- product_type: "Main" per il device, "Accessory" per cover/caricatore,
  "Part" per ricambi (batteria, schermo).
- IMPORTANT: Resolve pronouns using the CONVERSATION CONTEXT.
- Target Seller: extract if mentioned ("venduto da pegaso_italia").
- Output JSON only.

Query: {json.dumps(query, ensure_ascii=False)}
"""
\end{lstlisting}

\subsection{Validazione e Post-Elaborazione}
\label{subsec:parser-validation}

L'output LLM non viene usato direttamente ma passa attraverso due fasi
di validazione:

\begin{enumerate}
    \item \texttt{validate\_llm\_result()}: normalizza brand names,
        condizioni (``usato'' → \texttt{"used"}), prezzi (gestisce
        valori negativi o non numerici).
    \item \texttt{enforce\_numeric\_consistency()}: verifica che i range
        di prezzo siano coerenti (min < max) e che i valori numerici siano
        nel range ragionevole per il dominio e-commerce.
\end{enumerate}

La funzione \texttt{compute\_confidence()} assegna infine uno score
di confidenza 0–0.95 basato sui campi estratti con successo
(brand, constraints, product name).


% ─────────────────────────────────────────────────────────────
\section{Query Expansion e Named Entity Recognition}
\label{sec:query-expansion}
% ─────────────────────────────────────────────────────────────

\subsection{Query Expansion per il Retrieval}
\label{subsec:query-expansion}

Prima che la query raggiunga l'indice vettoriale Qdrant e il motore BM25,
viene arricchita semanticamente da un'apposita chiamata LLM in
\texttt{app/services/rag/query\_expansion.py}. Il prompt è volutamente
minimale e vincola l'output a massimo 10 parole:

\begin{lstlisting}[language=Python,caption={Prompt di query expansion (app/services/rag/query\_expansion.py)},label=lst:query-expansion]
prompt = f"""
You are an AI search assistant for an e-commerce platform.
Your task is to expand the user's short query into a slightly more
descriptive, keyword-rich query.
Add relevant brand names, categories, or synonyms if applicable.
Keep it under 10 words. Do NOT include phrases like
'here is the expansion' or quotes.
Just return the expanded query text.

User query: "{query}"
"""
\end{lstlisting}

Ad esempio, la query ``\emph{compara iphone 15 max}'' diventa
``\emph{Apple iPhone 15 Pro Max smartphone comparison}'', migliorando
il recall sia del retrieval denso che di quello sparso.

\subsection{Named Entity Recognition su Titoli Prodotto}
\label{subsec:ner}

Il modulo \texttt{app/services/nlp\_ner.py} usa un prompt NER per
estrarre attributi strutturati dai titoli dei prodotti eBay. Supporta
l'estrazione \textbf{batch} (più titoli in una singola chiamata):

\begin{lstlisting}[language=Python,caption={NER\_SYSTEM\_PROMPT (app/services/nlp\_ner.py)},label=lst:ner-prompt]
NER_SYSTEM_PROMPT = """
Sei un esperto di Information Retrieval per e-commerce.
Estrai attributi strutturati dal titolo di un prodotto eBay in JSON.
Devi estrarre:
- brand: il marchio (es. Apple, Samsung, Nike)
- model: il modello specifico (es. iPhone 13 Pro, Galaxy S21)
- specs: dizionario di caratteristiche tecniche
  (es. {"storage":"256GB","color":"Midnight"})
- category: categoria macro (es. "Smartphone", "Laptop")

IMPORTANTE:
1. Rispondi SOLO in formato JSON valido.
2. Non aggiungere spiegazioni o testo extra.
3. Se un attributo non e' presente, usa null.
4. Mantieni i valori delle specifiche concisi.
"""
\end{lstlisting}

Il risultato viene usato per arricchire i metadati dei documenti indicizzati
in Qdrant, migliorando la precisione del retrieval per query con filtri
su brand o specifiche tecniche.


% ─────────────────────────────────────────────────────────────
\section{Sentiment Analysis Ibrida}
\label{sec:sentiment-hybrid}
% ─────────────────────────────────────────────────────────────

Il modulo \texttt{app/services/nlp\_sentiment.py} implementa un sistema
di sentiment analysis a tre livelli che bilancia costo computazionale e
accuratezza. L'approccio è \textbf{lazy}: il livello LLM, il più costoso,
viene attivato solo come fallback finale.

\subsection{Tier 1 — Embedding con Anchor Vectors}
\label{subsec:sentiment-tier1}

Il livello primario usa \textbf{SentenceTransformer} per calcolare la
similarità coseno tra il testo da analizzare e un insieme di vettori
\emph{ancora} bilingui (italiano + inglese):

\begin{lstlisting}[language=Python,caption={Anchor vectors per sentiment embedding (app/services/nlp\_sentiment.py)},label=lst:sentiment-anchors]
positive_examples = [
    "excellent seller fast shipping",
    "great product highly recommended",
    "perfect transaction very happy",
    "ottimo venditore spedizione veloce",
    "perfetto consigliato"
]

negative_examples = [
    "terrible seller never again",
    "fake item scam seller",
    "bad service broken product",
    "venditore pessimo prodotto rotto",
    "truffa non comprate"
]
\end{lstlisting}

Il punteggio grezzo è derivato come differenza pesata tra similarità
media positiva e negativa, poi normalizzato tramite sigmoide logistica:
\[ \text{score} = \sigma(3 \cdot \text{raw}) = \frac{1}{1 + e^{-3 \cdot \text{raw}}} \]

\subsection{Tier 2 — Fallback Euristico}
\label{subsec:sentiment-tier2}

Se il modello di embedding non è disponibile, si attiva un contatore
lessicale su dizionari di parole positive e negative in italiano e inglese.
La formula applicata è analoga alla sigmoide del tier~1:
\[ \text{ratio} = \frac{\text{pos} - \text{neg}}{\text{tot} + 1} \]

\subsection{Tier 3 — Classificatore LLM}
\label{subsec:sentiment-tier3}

Come ultimo livello, un prompt LLM richiede una classificazione a
singola parola per testi ambigui:

\begin{lstlisting}[language=Python,caption={Prompt LLM per classificazione sentiment (app/services/nlp\_sentiment.py)},label=lst:sentiment-llm]
prompt = f"""
Analyze the following seller feedback comment and classify its
sentiment. Return ONLY ONE WORD from this list:
POSITIVE, NEGATIVE, NEUTRAL.
No other text, no punctuation.

Comment: "{text}"
"""
\end{lstlisting}

L'output a singola parola minimizza il rischio di parsing failure
e riduce la latenza della chiamata LLM.

\subsection{Sommario AI del Venditore}
\label{subsec:seller-summary}

La pipeline venditore (\texttt{app/services/seller\_pipeline.py}) usa
un prompt aggiuntivo per generare un sommario narrativo in italiano
dei feedback raccolti, limitando rigidamente l'output a 50 parole:

\begin{lstlisting}[language=Python,caption={Prompt sommario venditore (app/services/seller\_pipeline.py)},label=lst:seller-summary]
prompt = (
    f"Analizza queste recenti e autentiche recensioni di acquirenti "
    f"sul venditore '{seller_name}'. "
    "Scrivi in lingua italiana al massimo 2 o 3 frasi (max 50 parole) "
    "che riassumono il giudizio generale, i pregi e i potenziali "
    "difetti o rischi emersi. Sii conciso e diretto, senza convenevoli."
    "\n\n"
) + " ; ".join(texts)
\end{lstlisting}


% ─────────────────────────────────────────────────────────────
\section{Vision Prompting Multimodale}
\label{sec:vision-prompting}
% ─────────────────────────────────────────────────────────────

Quando l'utente allega un'immagine alla propria ricerca, il sistema
attiva una pipeline di \textbf{visual query enrichment} che opera
\emph{prima} del ciclo ReAct: l'immagine viene analizzata da un modello
multimodale e il risultato viene fuso con la query testuale originale.
L'intero flusso è implementato in \texttt{app/llm/vision.py} e
orchestrato da \texttt{app/agent/ebay\_agent.py}.

\subsection{Pipeline end-to-end}
\label{subsec:vision-pipeline}

\begin{figure}[htbp]
    \centering
    \fbox{\parbox{0.88\linewidth}{\centering\vspace{1.8cm}
        \textit{[Diagramma: immagine (base64) → Qwen-VL → JSON →
        query enrichment → loop ReAct]}\\
        \vspace{1.8cm}}}
    \caption{Pipeline di Vision Prompting nel sistema MCP\_ECOM.}
    \label{fig:vision-pipeline}
\end{figure}

Il flusso si articola in quattro fasi:

\begin{enumerate}
    \item \textbf{Decodifica base64}: l'immagine allegata dall'utente
        (campo \texttt{request.image}) viene decodificata da base64
        a bytes, con rimozione automatica dell'header data-URI
        (\texttt{data:image/png;base64,...}).

    \item \textbf{Chiamata Qwen-VL}: i bytes vengono inviati al modello
        \texttt{OLLAMA\_VISION\_MODEL} tramite \texttt{call\_ollama\_cloud()},
        con il \texttt{VISION\_PROMPT} come istruzione e l'immagine
        nel campo \texttt{images}.

    \item \textbf{Estrazione JSON}: la risposta grezza viene prima
        ripulita dai backtick markdown (regex \texttt{```json...```}),
        poi i delimitatori \texttt{\{...\}} vengono individuati per
        estrarre il JSON valido. In caso di fallimento del parsing,
        il testo grezzo viene avvolto in un dict di fallback con
        \texttt{confidence: 0.5}.

    \item \textbf{Enrichment della query}: brand e i primi 4 tag vengono
        concatenati alla query testuale originale. Se la query è vuota,
        viene costruita da zero come \texttt{"Trova: <brand> <tag1> ..."}.
\end{enumerate}

\subsection{Il Prompt Vision}
\label{subsec:vision-prompt}

Il prompt istruisce il modello con un ruolo da \emph{esperto e-commerce}
e impone un output JSON con schema fisso a 5 campi:

\begin{lstlisting}[language=Python,caption={VISION\_PROMPT (app/llm/vision.py)},label=lst:vision-prompt]
VISION_PROMPT = """Analizza questa immagine come un esperto e-commerce
per una ricerca prodotto. Estrai i dettagli e restituisci
ESCLUSIVAMENTE un resoconto in formato JSON.
{
  "description": "Descrizione estensiva in italiano (marca, modello,
                  colore, materiale, dettagli identificativi)",
  "tags": ["keyword1", "keyword2", "keyword3"],
  "brand": "Nome Brand se chiaramente visibile (oppure null)",
  "condition_clues": "Condizioni apparenti in italiano oppure null",
  "confidence": 0.95
}
NON inserire spiegazioni, introduzioni, markdown o backtick.
Produci SOLO output JSON valido.
"""
\end{lstlisting}

La regola ``JSON only'' è fondamentale perché l'output viene parsato
programmaticamente: qualsiasi testo aggiuntivo causerebbe un fallback
al dict minimale con \texttt{confidence: 0.5}.

\subsection{Query Enrichment e Trasmissione al Planner}
\label{subsec:vision-integration}

Nell'agente, la query viene arricchita in-place prima che il loop
ReAct parta. Il brand estratto e i tag (max 4) vengono accodati
tra parentesi:

\begin{lstlisting}[language=Python,caption={Enrichment della query e VisionAnalysisEvent (app/agent/ebay\_agent.py)},label=lst:vision-enrichment]
# Costruisce il supplemento testuale dalla visione
search_additive = []
if brand:
    search_additive.append(brand)
if tags:
    search_additive.extend(tags[:4])       # max 4 tag

compact_str = " ".join(search_additive)

# Fonde con la query originale
if not original_query:
    request.query = f"Trova: {compact_str}"
else:
    request.query = f"{original_query} ({compact_str})"

# Emette un evento SSE dedicato per la card visione nell'UI
yield VisionAnalysisEvent(
    description=desc_text,
    tags=tags,
    brand=brand,
    condition_clues=condition,
    confidence=confidence
).model_dump()
\end{lstlisting}

Parallelamente, la \texttt{description} completa viene salvata nello
scratchpad (\texttt{vision\_description}) e iniettata nel prompt del
planner ad ogni step come riga di contesto aggiuntiva:

\begin{lstlisting}[language=Python,caption={Iniezione nel prompt del planner (app/agent/prompts.py)},label=lst:vision-injection]
vision_desc = scratchpad.get("vision_description")
vision_context = f"Vision description: {vision_desc}\n" \
                 if vision_desc else ""
\end{lstlisting}

In questo modo il planner dispone sia della query arricchita (canale
principale per i tool) sia della descrizione narrativa completa
(contesto semantico per il ragionamento). Ad esempio, la foto di un
paio di scarpe Nike rosse trasforma la query ``\emph{cerca online}''
in ``\emph{cerca online (Nike Air Max red sneakers leather)}'' e
fornisce al modello il dettaglio completo sulle condizioni visibili
dell'articolo.


% ─────────────────────────────────────────────────────────────
\section{Context Compaction e Budget Management}
\label{sec:context-compaction}
% ─────────────────────────────────────────────────────────────

Per contenere le dimensioni dei prompt entro il budget di 6000 token,
il sistema applica tre funzioni di compattazione che filtrano i dati
rilevanti da strutture potenzialmente molto grandi.

\subsection{Compattazione del Catalogo Strumenti}
\label{subsec:compact-tools}

La funzione \texttt{\_compact\_tool\_catalog\_for\_prompt()} riduce ogni
voce del catalogo MCP a soli 10 campi essenziali, escludendo payload
pesanti come i metadati di configurazione avanzati:

\begin{center}
\begin{tabular}{ll}
\hline
\textbf{Campo incluso} & \textbf{Scopo nel prompt} \\
\hline
\texttt{description} & Descrive il tool all'LLM \\
\texttt{input\_schema} & Schema dei parametri accettati \\
\texttt{required\_fields} & Parametri obbligatori \\
\texttt{tags} & Etichette semantiche \\
\texttt{examples} & Max 2 esempi d'uso \\
\texttt{state\_key} & Dove archiviare il risultato \\
\texttt{cost} & Indicatore di costo relativo \\
\texttt{latency\_class} & Veloce / lento / very-slow \\
\texttt{dependencies} & Tool da invocare prima \\
\texttt{produced\_entities} & Entità generate (product, seller, …) \\
\hline
\end{tabular}
\end{center}

\subsection{Compattazione dello Scratchpad}
\label{subsec:compact-scratchpad}

Lo scratchpad dell'agente può contenere decine di campi. La funzione
\texttt{\_compact\_scratchpad\_for\_prompt()} seleziona i 20+ campi
rilevanti per il ragionamento del planner, escludendo payload di dati
grezzi (i risultati completi di ricerca, i chunk RAG, ecc.).

\subsection{Compattazione dei Dati Finali}
\label{subsec:compact-final-data}

Prima della fase di final answer, \texttt{\_compact\_final\_data\_for\_prompt()}
costruisce una vista compatta che include:

\begin{itemize}
    \item I primi 6 risultati di ricerca (\texttt{top\_results[:6]})
    \item Un record venditore ridotto (name, count, trust, sentiment)
    \item Dati di confronto e metriche IR
    \item I campi \texttt{data} e \texttt{results} dei \texttt{tool\_states}
          vengono esplicitamente rimossi per evitare payload oltre i 400 KB
          che causerebbero errori HTTP~400 da Ollama
\end{itemize}


% ─────────────────────────────────────────────────────────────
\section{Riepilogo delle Strategie}
\label{sec:prompt-summary}
% ─────────────────────────────────────────────────────────────

La tabella seguente riassume tutte le strategie di prompting adottate
nel progetto MCP\_ECOM, con i file di riferimento e il numero di
occorrenze:

\begin{table}[htbp]
\centering
\small
\begin{tabular}{p{4cm} p{4.5cm} p{5cm}}
\hline
\textbf{Strategia} & \textbf{File} & \textbf{Applicazione} \\
\hline
Role Prompting & \texttt{prompts.py}, \texttt{conversation.py} &
    4 personalità distinte (planner, playwright, final, chat) \\
\hline
Chain-of-Thought & \texttt{prompts.py}, \texttt{planner.py} &
    Campo \texttt{thought} in italiano + intent esplicito \\
\hline
Structured Output (JSON) & \texttt{parser.py}, \texttt{vision.py},
    \texttt{ltr\_trainer.py} &
    Schema con validazione post-LLM \\
\hline
Tone Control dinamico & \texttt{prompts.py}, \texttt{ebay\_agent.py} &
    3 livelli (amichevole / professionale / neutral) + override \\
\hline
Context Injection & \texttt{ebay\_agent.py}, \texttt{prompts.py} &
    Profilo utente MCP, preferenze LTM, storico sessione \\
\hline
Pronoun Resolution & \texttt{parser.py} &
    CONVERSATION CONTEXT block per pronomi impliciti \\
\hline
Query Expansion & \texttt{query\_expansion.py} &
    Arricchimento semantico pre-retrieval (max 10 parole) \\
\hline
Named Entity Recognition & \texttt{nlp\_ner.py} &
    Estrazione batch brand/model/specs da titoli prodotto \\
\hline
LTR Judge Prompt & \texttt{ltr\_trainer.py} &
    Scala 0–4 con regola vincoli (max~1 se violazione) \\
\hline
Sentiment Ibrido & \texttt{nlp\_sentiment.py} &
    Embedding → euristico → LLM (tier progressivo) \\
\hline
Vision Prompting & \texttt{vision.py} &
    Qwen-VL image-to-JSON, iniettato nella query \\
\hline
Token Budgeting & \texttt{prompts.py} &
    6000 token max, troncamento selettivo per campo \\
\hline
Priority Injection & \texttt{ebay\_agent.py}, \texttt{prompts.py} &
    custom $>$ tone $>$ base $>$ context \\
\hline
Filtri Negativi & \texttt{prompts.py} &
    Esclusione esplicita attributi vietati nella risposta finale \\
\hline
\end{tabular}
\caption{Riepilogo delle strategie di prompting nel progetto MCP\_ECOM.}
\label{tab:prompt-strategies}
\end{table}

\begin{figure}[htbp]
    \centering
    \fbox{\parbox{0.85\linewidth}{\centering\vspace{2.5cm}
        \textit{[Diagramma: flusso di composizione del prompt —
        custom instructions → tone → system prompt base → context
        → scratchpad compatto → query utente]}\\
        \vspace{2.5cm}}}
    \caption{Pipeline di composizione del prompt nel sistema MCP\_ECOM.}
    \label{fig:prompt-pipeline}
\end{figure}

In sintesi, il progetto MCP\_ECOM tratta il prompting come un artefatto
ingegneristico di prima classe: ogni chiamata LLM è governata da istruzioni
costruite programmaticamente, validate e adattate al contesto corrente.
La combinazione di role prompting, chain-of-thought, output strutturato,
tone control dinamico e context injection a più livelli consente al sistema
di produrre risposte coerenti, personalizzate e affidabili attraverso
l'intero ciclo di vita della richiesta e-commerce.
